\documentclass[12pt,letterpaper]{article}

\usepackage{amsmath}
%\usepackage[dvips]{graphicx}
%\usepackage{epic}

% revise margins
\setlength{\headheight}{0in}
\setlength{\topmargin}{0in}
\setlength{\headsep}{0in}
\setlength{\textheight}{9.0in}
\setlength{\footskip}{0.4in}
\setlength{\oddsidemargin}{0in}
\setlength{\evensidemargin}{0in}
\setlength{\parindent}{0in}
\setlength{\parsep}{12pt}
\setlength{\textwidth}{6.5in}

% san serif font
\renewcommand{\familydefault}{cmss}

% 3rd itemize level is + rather than *
\renewcommand{\labelitemiii}{$\circ$}

\newcommand{\E}{\text{E}}       % expected value
\newcommand{\var}{\text{var}}   % variance
\newcommand{\cov}{\text{cov}}   % covariance
\newcommand{\cor}{\text{cor}}   % correlation
\renewcommand{\Pr}{\text{Pr}}     % probability
\newcommand{\iid}{\text{ iid }} % iid

% spacing commands
\newcommand{\HRule}{\rule{\linewidth}{4pt}}
\newcommand{\hsphalf}{\hspace*{0.5cm}}
\newcommand{\hspone}{\hspace*{1cm}}
\newcommand{\hsptwo}{\hspace*{2cm}}
\newcommand{\vsphalf}{\vspace*{0.5cm}}
\newcommand{\vspone}{\vspace*{1cm}}
\newcommand{\vsptwo}{\vspace*{2cm}}

% page number centered at bottom
\pagestyle{plain}

\begin{document}

\fontsize{30}{30} \selectfont

\centerline{\textbf{Numerical integration}}

\HRule

\vsphalf

\fontsize{20}{25} \selectfont

To be fancy, our topic is \emph{quadrature}.

\bigskip

We seek $I = \int_a^b f(x) \; dx$.

\bigskip \bigskip

\textbf{Method 1}: \emph{The extended midpoint rule}
\begin{enumerate}
\item Pick $N$ large.
\item Let $x_i = a + (i - 1/2)h$, $h = (b-a)/N$, \\
   for $i = 1, \ldots, N$.
\item Let $f_i = f(x_i)$
\item Then $I \approx h \sum_i f_i$
\end{enumerate}

\bigskip

\textbf{Sample code}:
\begin{verbatim}
    emr <-
    function(f,a,b,n=1000)
    {
      h <- (b-a)/n
      h*sum(f( seq(a+h/2, b, by=h) ))
    }
\end{verbatim}

\textbf{Points}:
\begin{itemize}
\item This is the simplest thing to do
\item Error $O(1/N^2)$
\item Don't use this in production
\end{itemize}


\newpage

\fontsize{30}{30} \selectfont

\centerline{\textbf{The basics}}

\HRule

\vsphalf

\fontsize{20}{25} \selectfont

In the following, we consider $x_0, x_1, \ldots x_N$ such that
$x_i = x_0 + i h$, and write $f_i = f(x_i)$.

\bigskip

\textbf{Note}: The common choice is $x_0 = a$ and $h = (b-a)/N$,

\hspone so that $x_N = b$.

\bigskip \bigskip

\textbf{Closed formula}: use $x_0, x_N$

\bigskip

\textbf{Open formula}: don't use $x_0, x_N$ (e.g., when it's hard
to calculate $f$ at one of $a$ or $b$)

\bigskip \bigskip

\textbf{Basic rules}

\bigskip

\hspone \emph{Trapezoidal rule}: \hfill $\int_{x_0}^{x_1} f(x) \; dx
\approx h \{ \frac{1}{2} f_0 + \frac{1}{2} f_1 \}$

\bigskip

\hspone \emph{Simpson's rule}: \hfill $\int_{x_0}^{x_2} f(x) \: dx
\approx h \{ \frac{1}{3} f_0 + \frac{4}{3} f_1 + \frac{1}{3} f_2 \}$

\bigskip \bigskip

\hspone \emph{Extrapolation formulas}:
\hfill $\int_{x_0}^{x_1} f(x) \: dx \approx h \{ f_1 \}$

\bigskip

\hfill $\int_{x_0}^{x_1} f(x) \: dx \approx h \{ \frac{3}{2} f_1 -
\frac{1}{2} f_2 \}$

\bigskip

\hfill $\int_{x_0}^{x_1} f(x) \: dx \approx h \{ \frac{23}{12} f_1 -
\frac{16}{12} f_2 + \frac{5}{12}f_3 \}$

\newpage

\fontsize{30}{30} \selectfont

\centerline{\textbf{The basics} (continued)}

\HRule

\vsphalf

\fontsize{20}{25} \selectfont

\textbf{Extended trapezoidal rule (closed)}
$$ \int_{a}^{b} f(x) \; dx = h \{ {\textstyle \frac{1}{2}} f_0 + f_1 + f_2 + \ldots +
f_{N-1} + {\textstyle \frac{1}{2}} f_N \} + O({\textstyle \frac{1}{N^2}})$$

\bigskip \bigskip

\textbf{Extended Simpson's rule (closed)}
\begin{eqnarray*}
\int_{a}^{b} f(x) \; dx & = & h \{ {\textstyle \frac{1}{3}}f_0 +
{\textstyle \frac{4}{3}}f_1 + {\textstyle \frac{2}{3}}f_2 + {\textstyle \frac{4}{3}}f_3
+ \ldots \\
& & \hsptwo + {\textstyle \frac{2}{3}}f_{N-2} + {\textstyle
  \frac{4}{3}}f_{N-1} + {\textstyle \frac{1}{3}}f_N \} + O({\textstyle
  \frac{1}{N^4}})
\end{eqnarray*}

\vsptwo

\textbf{Extended trapezoidal rule (open)}
$$ \int_{a}^{b} f(x) \; dx = h \{ {\textstyle \frac{3}{2}} f_1 + f_2 + \ldots +
f_{N-2} + {\textstyle \frac{3}{2}} f_{N-1} \} + O({\textstyle \frac{1}{N^2}})$$

\bigskip \bigskip

\textbf{Extended Simpson's rule (open)}
\begin{eqnarray*}
\int_{a}^{b} f(x) \; dx & = &h \{ {\textstyle \frac{27}{12}} f_1 + 0 +
  {\textstyle \frac{13}{12}} f_3 + {\textstyle \frac{4}{3}} f_4 +
{\textstyle \frac{2}{3}} f_5 + \ldots \\
& & + {\textstyle \frac{4}{3}} f_{N-4} + {\textstyle \frac{13}{12}}
  f_{N-3} + 0 + {\textstyle \frac{27}{12}} f_{N-1} \} + O({\textstyle \frac{1}{N^4}})
\end{eqnarray*}

\newpage

\fontsize{30}{30} \selectfont

\centerline{\textbf{Trapezoidal rule: Implementation}}

\HRule

\vsphalf

\fontsize{20}{25} \selectfont

An important point regarding the trapezoidal rule: to go from $N$ to
$2N$, you can make use of your previous work.

\bigskip

We make use of the following basic engine for the trapezoidal rule.  I
assume that the function \verb|f| is vectorized and that \verb|n| is a
power of 2.

\fontsize{18}{18} \selectfont

\begin{verbatim}
    subtrap <-
    function(f,a,b,n=1)
    {
      h <- (b-a)/n
      if(n==1) return( h*mean( f(c(a,b)) ) )
      else return(h*sum( f(seq(a+h,b,by=2*h)) ))
    }
\end{verbatim}

\fontsize{20}{25} \selectfont

We use the above function by taking $I_{i+1} = \frac{1}{2} I_i +
S_{i+1}$, where $S_i$ is the output from the function \verb|subtrap|
and $I_i$ is the estimated integral using $N = 2^{i-1}$ steps.

\fontsize{18}{18} \selectfont

\begin{verbatim}
    trap <-
    function(f,a,b,tol=1e-8,maxit=1000)
    {
      i.old <- subtrap(f,a,b,1); n <- 2
      for(i in 2:maxit) {
        s <- subtrap(f,a,b,n)
        i.new <- i.old/2 + s
        if(abs(i.new-i.old) < tol) break
        i.old <- i.new
        n <- n*2
      }
      i.new
    }
\end{verbatim}

\newpage

\fontsize{30}{30} \selectfont

\centerline{\textbf{Simpson's rule: Implementation}}

\HRule

\vsphalf

\fontsize{20}{25} \selectfont

It is interesting to note that you can use the \verb|subtrap|
function to obtain Simpson's rule.

\bigskip

Let $S_i$ be the output from \verb|subtrap| with $N = 2^{i-1}$, and
let $I_i$ be our estimate of the integral at step $i$.

\bigskip

For Simpson's rule, we take $I_1 = 2S_1/3$ and then $I_{i+1} = I_i/2 +
(4S_{i+1} - S_i)/3$.

\fontsize{18}{18} \selectfont

\begin{verbatim}
    simp <-
    function(f,a,b,tol=1e-8,maxit=1000)
    {
      i.old <- subtrap(f,a,b,1)*2/3
      n <- 2;  old.s <- 0
      for(i in 2:maxit) {
        s <- subtrap(f,a,b,n)
        i.new <- (i.old/2 + (4*s - old.s)/3)
        if(abs(i.new-i.old) < tol) break
        i.old <- i.new; old.s <- s
        n <- n*2
      }
      i.new
    }
\end{verbatim}

\fontsize{20}{25} \selectfont

Consider the example $\int_0^2 x^2(1-x)sin(x^2) \; dx$.

\bigskip

With \verb|tol=1e-8|, Simpson's rule uses $N = 512$ while the
trapezoidal rule uses $N = 512 \times 128$.

\newpage

\fontsize{30}{30} \selectfont

\centerline{\textbf{Further explanation}}

\HRule

\vsphalf

\fontsize{20}{25} \selectfont

Let $S(h)$ and $T(h)$ denote the approximations from Simpson's rule
and the trapezoidal rule, respectively.

\bigskip

Then $S(h) = \frac{1}{3} \{ 4 T(h) - T(2h) \}$

\bigskip \bigskip

The \verb|subtrap| function provides $X(h) = T(h) - \frac{1}{2}
T(2h)$.

\begin{eqnarray*}
\text{Thus } S(h) & = & {\textstyle \frac{1}{3}} \{ 4 T(h) - T(2h) \}
\\ \\
& = & {\textstyle \frac{1}{3}} \{ [2 T(2h) - {\textstyle \frac{1}{2}}
            T(4h)] + [4 T(h) - \\
            & & \hspone \hsphalf 2 T(2h)] - [T(2h) + {\textstyle \frac{1}{2}}
T(4h)] \} \\ \\
& = & {\textstyle \frac{1}{2}} \cdot {\textstyle \frac{1}{3}} \{ 4
            T(2h) - T(4h) \} + {\textstyle \frac{4}{3}} \{T(h) - \\
            & & \hspone \hsphalf {\textstyle \frac{1}{2}} T(2h) \} -
            {\textstyle \frac{1}{3}} \{ T(2h) - {\textstyle \frac{1}{2}}
            T(4h) \} \\ \\
& = & {\textstyle \frac{1}{2}}S(2h) + {\textstyle \frac{1}{3}} \{
            4 X(h) - X(2h) \}
\end{eqnarray*}


\newpage

\fontsize{30}{30} \selectfont

\centerline{\textbf{Romberg's algorithm}}

\HRule

\vsphalf

\fontsize{20}{25} \selectfont

Simpson's rule combines $T(h)$ and $T(2h)$ in a way that some error
terms are cancelled.  In Romberg's algorithm, this method is
continued.

\bigskip

If the function being integrated is sufficiently smooth, impressive
gains are obtained.

\vspone

Let $T_{m0} = T[(b-a)2^{-m}]$ and define
\begin{eqnarray*}
T_{mn} & = & \frac{4^n T_{m,n-1} - T_{m-1,n-1}}{4^n-1}\\ \\
       & = & T_{m,n-1} - {\textstyle \frac{1}{4^n-1}} \{ T_{m-1,n-1} -
       T_{m,n-1} \}
\end{eqnarray*}

\vspone

The Romberg matrix:
$$ \begin{pmatrix} T_{00} \\
T_{10} & T_{11} \\
T_{20} & T_{21} & T_{22} \\
T_{30} & T_{31} & T_{32} & T_{33} \\
\vdots & \vdots & \vdots & \vdots \end{pmatrix} $$

\vspone

Recall that if $X_i$ is the value from the $i$th call to
\verb|subtrap|, $T_{00} = X_1$ and $T_{i0} = X_{i+1} + \frac{1}{2}
T_{i-1,0}$.

\newpage

\fontsize{30}{30} \selectfont

\centerline{\textbf{Improper integrals}}

\HRule

\vsphalf

\fontsize{20}{25} \selectfont

The above methods need to be revised in the case that
\begin{enumerate}
\item The integrand goes to a finite limit at finite upper and lower
limits, but cannot be calculated right on one of the limits (e.g.,
$\sin x/x$ at $x=0$).

\item The upper limit of integration is $\infty$ or the lower limit is
$-\infty$.

\item There is a singularity at either limit (e.g., $x^{-1/2}$ at
$x=0$.

\item There is a singularity at a known spot between the limits.
\end{enumerate}

\bigskip

\textbf{Solution}: Replace the function \verb|subtrap| with a version
of the extended midpoint rule, tripling $N$ at each call.

\setlength{\unitlength}{1.5cm}
\begin{center} \begin{picture}(10.5,3.0)

% lines
\put(0,0.5){\line(1,0){9}}
\put(0,0.35){\line(0,1){0.3}}
\put(9,0.35){\line(0,1){0.3}}
\put(0,1.5){\line(1,0){9}}
\put(0,1.35){\line(0,1){0.3}}
\put(9,1.35){\line(0,1){0.3}}
\put(0,2.5){\line(1,0){9}}
\put(0,2.35){\line(0,1){0.3}}
\put(9,2.35){\line(0,1){0.3}}

% internal ticks
\put(1,0.35){\line(0,1){0.3}}
\put(2,0.35){\line(0,1){0.3}}
\put(3,0.35){\line(0,1){0.3}}
\put(4,0.35){\line(0,1){0.3}}
\put(5,0.35){\line(0,1){0.3}}
\put(6,0.35){\line(0,1){0.3}}
\put(7,0.35){\line(0,1){0.3}}
\put(8,0.35){\line(0,1){0.3}}
\put(9,0.35){\line(0,1){0.3}}
\put(3,1.35){\line(0,1){0.3}}
\put(6,1.35){\line(0,1){0.3}}

% arrows
\thicklines
\put(4.5,2.9){\vector(0,-1){0.4}}
\put(4.5,1.1){\vector(0,1){0.4}}
\put(4.5,0.1){\vector(0,1){0.4}}

\put(1.5,1.9){\vector(0,-1){0.4}}
\put(1.5,0.1){\vector(0,1){0.4}}
\put(7.5,1.9){\vector(0,-1){0.4}}
\put(7.5,0.1){\vector(0,1){0.4}}

\put(0.5,0.9){\vector(0,-1){0.4}}
\put(2.5,0.9){\vector(0,-1){0.4}}
\put(3.5,0.9){\vector(0,-1){0.4}}
\put(5.5,0.9){\vector(0,-1){0.4}}
\put(6.5,0.9){\vector(0,-1){0.4}}
\put(8.5,0.9){\vector(0,-1){0.4}}

%text
\put(9.5,0.35){$N = 3^2$}
\put(9.5,1.35){$N = 3^1$}
\put(9.5,2.35){$N = 3^0$}
\end{picture} \end{center}

Use code like

{\Large \verb|  x <- c(seq(a+h/2,b,by=3*h), seq(a+2.5*h,b,by=3*h))|}

\smallskip

We can easily revise \verb|trap| and \verb|simp| to use this.  We thus
solve issue 1.




\newpage

\fontsize{30}{30} \selectfont

\centerline{\textbf{Improper integrals} (continued)}

\HRule

\vsphalf

\fontsize{20}{25} \selectfont

To solve the other issues, we use two techniques:
\begin{enumerate}
\item Change of variables.
\item Break up the integral into pieces
\end{enumerate}

\bigskip

For example, if $a > 0$ and $f(t) \to 0$ faster than $1/t^2 \to 0$ as
$t \to \infty$, then we can use $u = 1/t$ as follows:

\bigskip

\hspone $\int_a^\infty f(t) \; dt = \int_{0}^{1/a} \frac{1}{u^2}
f(\frac{1}{u}) \; du$

\bigskip

This also works if $b < 0$ and the lower limit is $-\infty$.

\bigskip

I think that's enough on that; you can see the general direction.  See
Lange \S *.* for additional advice and examples.





\newpage

\fontsize{30}{30} \selectfont

\centerline{\textbf{Gaussian quadrature}}

\HRule

\vsphalf

\fontsize{20}{25} \selectfont

The amazingly cool technique of \emph{Gaussian quadrature\/} is
appropriate when we wish to integrate $f(x)$ against a probability
measure $\mu$.

\bigskip

For fixed points $x_i$ and positive weights $w_i$,
$$ \int_{-\infty}^{\infty} f(x) \; d\mu(x) \approx \sum_{i=0}^{k}
w_i f(x_i)$$

\bigskip

For the trapezoidal rule, $\mu = \text{uniform}(a,b)$ and the $x_i$
are uniformly spaced.  Generally, $\mu$ is non-uniform and the $x_i$
cluster in regions of high probability.

\bigskip

If $\mu$ possesses an orthonormal polynomial sequence $\psi_n(x)$ for
$n=0, 1, 2, \ldots$, then $x_0, x_1, \ldots, x_k$ are the roots of
$\psi_{k+1}(x)$.

\vspone

\textbf{Fact}: If $\mu$ is not concentrated at a finite number of
points, then the roots $\{x_i\}$ are real and distinct, and there
exist positive weights $w_i$ such that the quadrature formula above is
exact when $f(x)$ is a polynomial of degree $\le 2k+1$.


\newpage

\fontsize{30}{30} \selectfont

\centerline{\textbf{Orthonomial polynomials} (continued)}

\HRule

\vsphalf

\fontsize{20}{25} \selectfont

Consider the inner product
$$<f,g> = \int_{-\infty}^{\infty} f(x) g(x)
\; d\mu(x)$$

\bigskip

The orthonomial polynomials mentioned above satisfy

$<p_i,p_j> = \delta_{ij}$ (i.e., $= 0$ if $i \ne j$ and $=1$ if $i=j$)

\vspone

They can be constructed as follows:

\bigskip

Take $p_{-1}(x) = 0$ and $p_0(x) = 1$.

\bigskip

For $j \ge 0$, let $q_{j+1} = (x-a_j)p_j(x) - b_j p_{j-1}(x)$

\hspone and let $p_j = q_j / \sqrt{<q_j, q_j>}$

\vspone

My recommendation: if you come across an integral of the form $\int
f(x) \; d\mu(x)$, consider using Gaussian quadrature, and check
out \emph{Numerical Recipes in C}, \S *.*.

\bigskip

Rather than calculating the abscissas $x_i$ and weights $w_i$,
you may wish to pull them from tables (in an old and dusty book).

\newpage

\fontsize{30}{30} \selectfont

\centerline{\textbf{Example 1}}

\HRule

\vsphalf

\fontsize{20}{25} \selectfont

[Chakravarti et al., Genetics 128:175--182, 1991]

\bigskip

Suppose a genome consists of $k$ chromosomes of common length $L$, and
that we have typed $m$ markers on a set of fully informative
offspring.  Let $G = kL$ be the genome length.

\bigskip

Let $n_{ij}$ be the number of meioses typed on both markers $i$ and
$j$ and $r_{ij}$ be the number of recombinants observed.

\bigskip

We assume (falsely) no interference and (even more falsely) that the
pairs are independent, and wish to estimate $L$ (and hence $G$).

\bigskip

\textbf{Log likelihood}:

\begin{center} \begin{tabular}{l}
$l(L) = \sum_{i \ne j} \log \{ \int_{0}^{\theta'} \theta^{r_{ij}}
(1-\theta)^{n_{ij}-r_{ij}} f(\theta; L) \; d\theta $ \\
\multicolumn{1}{r}{$+ (k-1)(1/2)^{n_{ij}} \}$}
\end{tabular} \end{center}

where $\theta' = (1-e^{-2L})/2$ (ie, the maximum recombination
fraction for a chromosome of length $L$) and $f(\theta; L) = \{ 2L +
\log(1-2\theta)\}/\{L^2(1-2\theta)\}$.

\bigskip

We need to calculate that integral (and then maximize $l$ with respect
to $L$, to get the MLE $\hat{L}$).

\bigskip

\textbf{Major problem}: The integral can be \emph{really} small, hence
we have a problem with \emph{underflow}.


\newpage

\fontsize{30}{30} \selectfont

\centerline{\textbf{Example 1} (continued)}

\HRule

\vsphalf

\fontsize{20}{25} \selectfont

\textbf{The trick}:

\bigskip

Let $g_i = \log f_i$. Suppose we wish to calculate $\log(f_1 +
f_2)$.
\begin{eqnarray*}
\text{Then } \log(f_1 + f_2) & = & \log(e^{g_1} + e^{g_2}) \\
& = & \log\{ e^{g_1} ( 1 + e^{g_2 - g_1}) \} \\
& = & g_1 + \log(1 + e^{g_2 - g_1})
\end{eqnarray*}

\bigskip

``Obviously,'' a problem occurs when $g_2 \gg g_1$, in which case
$\log(f_1 + f_2) \approx g_2$, but the above formula will result in an
overflow.  Thus, I recommend the following:

\fontsize{20}{20}

\begin{verbatim}
    addlog <-
    function(a, b, thresh=200)
    {
      if(b > a + thresh) return(b)
      else if(a > b + thresh) return(a)
      else return(a + log(1+exp(b-a)))
    }
\end{verbatim}

\fontsize{20}{25}

Note that one may also want a function \verb|subtractlog()|.

\bigskip

We can modify any code to do numerical integration, to calculate
$\log\{\int_a^b f(x) \; dx\}$ working only with $\log f(x)$.  \\ This
will avoid the underflow problem.

\bigskip

To calculate $S = \log \sum_{i=1}^n f_i$, use

\hspone $S_1 = g_1$, $S_{i+1} = \verb|addlog|(S_i, g_{i+1})$

\newpage

\fontsize{30}{30} \selectfont

\centerline{\textbf{Example 2}}

\HRule

\vsphalf

\fontsize{20}{25} \selectfont

[Broman and Weber, Am J Hum Genet 66:1911--1926, 2000]

\bigskip

Suppose we observe vectors $y_1, y_2, \ldots, y_N$, the realizations
of a \emph{thinned} stationary gamma renewal process on the interval
$[0, L]$ (where $L$ is known), with shape parameter $\nu$.  We seek to
estimate $\nu$.

\bigskip

Let's focus on a single realization $y = (y_1, \ldots, y_m)$.

\bigskip

\setlength{\unitlength}{0.5cm}
\begin{center} \begin{picture}(30,2)

% lines
\put(0,0.5){\line(1,0){30}}
\put(0,0){\line(0,1){1}}
\put(30,0){\line(0,1){1}}

% Xs
\thicklines
\put(3.5,1){\line(1,-1){1}}
\put(4.5,1){\line(-1,-1){1}}
\put(7.5,1){\line(1,-1){1}}
\put(8.5,1){\line(-1,-1){1}}
\put(9.5,1){\line(1,-1){1}}
\put(10.5,1){\line(-1,-1){1}}
\put(16.5,1){\line(1,-1){1}}
\put(17.5,1){\line(-1,-1){1}}
\put(19.1,1){\line(1,-1){1}}
\put(20.1,1){\line(-1,-1){1}}
\put(23.3,1){\line(1,-1){1}}
\put(24.3,1){\line(-1,-1){1}}

% x's
\put(1.5,1.5){$x_0$}
\put(5.5,1.5){$x_1$}
\put(8.5,1.5){$x_2$}
\put(13,1.5){$x_3$}
\put(17.8,1.5){$x_4$}
\put(21.2,1.5){$x_5$}
\put(26.4,1.5){$x_n$}

\end{picture}

\bigskip \bigskip \bigskip

\begin{picture}(30,2)

% lines
\put(0,0.5){\line(1,0){30}}
\put(0,0){\line(0,1){1}}
\put(30,0){\line(0,1){1}}

% Xs
\thicklines
\put(3.5,1){\line(1,-1){1}}
\put(4.5,1){\line(-1,-1){1}}
\put(7.5,1){\line(1,-1){1}}
\put(8.5,1){\line(-1,-1){1}}
\put(16.5,1){\line(1,-1){1}}
\put(17.5,1){\line(-1,-1){1}}

% y's
\put(1.5,1.5){$y_0$}
\put(5.5,1.5){$y_1$}
\put(12.0,1.5){$y_2$}
\put(23.0,1.5){$y_m$}
\end{picture} \end{center}

\bigskip

Our model is that the points for the $X$ process come from a
stationary gamma renewal model with shape parameter $\nu$ and mean
inter-point distance $1/2$, so that
$$x_1, x_2, \ldots \sim \iid \text{gamma}(\nu, 2\nu)$$
with density $f(x; \nu) = e^{-2 \nu x}(2\nu)^\nu x^{\nu-1} /
\Gamma(\nu)$.

\bigskip

The distribution of $x_0$ is that which is required to give
stationarity: $g(x;\nu) = 2[1-F(x;\nu)]$ where $F$ is the cdf of $f$.
The last point $x_n$ is treated as censored.

\bigskip

\newpage

\fontsize{30}{30} \selectfont

\centerline{\textbf{Example 2} (continued)}

\HRule

\vsphalf

\fontsize{20}{25} \selectfont

The $Y$ process (what we observe) is obtained by thinning the $X$
process: we flip a coin at each point in the $X$ process,
independently; H $\to$ retain point; T $\to$ drop point.

\bigskip

The $Y$ process is also a renewal process.
$$y_1, y_2, \ldots \sim \iid f^\star(x;\nu) = \sum_{k\ge 1} (1/2)^2
f_k(y; \nu)$$
where $f_k$ is the density of a gamma distribution with shape and rate
parameters $k\nu$ and $2\nu$---the convolution of $f(y;\nu)$ with
itself $k$ times.

\bigskip

The distribution of $y_0$ is
\begin{eqnarray*}
g^\star(y;\nu) & = & (1/2) g(y;\nu) + \sum_{k \ge 1} (1/2)^{k+1} (g *
f_k)(y;\nu) \\
& = & 1 - F^\star(y;\nu)
\end{eqnarray*}
where $F^\star$ is the cdf of $f^\star$.  Let $G^\star$ denote the cdf
of $g^\star$.

\bigskip

Then $y$ contributes the following to the log likelihood for $\nu$:
$$l(\nu;y) = \left\{ \begin{tabular}{lc}
$\log[1 - G^\star(L;\nu)]$ & if $m = 0$ \\
$\log g^\star(y_0;\nu) + \log g^\star(y_1;\nu)$ & if $m=1$ \\
$\log g^\star(y_0;\nu) + \sum_{j=1}^{m-1} \log f^\star(y_j; \nu)$ & \\
\hsptwo $+ \log g^\star(y_m; \nu)$  & otherwise

\end{tabular} \right. $$

\newpage

\fontsize{30}{30} \selectfont

\centerline{\textbf{Example 2} (continued)}

\HRule

\vsphalf

\fontsize{20}{25} \selectfont

We seek to maximize $\sum_i l(\nu; y_i)$ where $y_i$ is a
vector like that we played with above.

\bigskip

To calculate $l(\nu; y_i)$, we need to be able to calculate
\begin{eqnarray*}
f^\star(y;\nu) & = & \sum_{k=1}^\infty (1/2)^k f_k(y;\nu) \\
g^\star(y;\nu) & = & 1 - F^\star(y; \nu) = \int_{y}^\infty
f^\star(t;\nu) \; dt \\
1 - G^\star(y;\nu) & = & \int_{y}^\infty [1-F^\star(t;\nu)] \; dt
\end{eqnarray*}

\bigskip

$f^\star(y;\nu)$ was calculated using (``by use of'') its explicit
formula, summing over $k$ from 1 to 25.

\bigskip

$g^\star(y;\nu)$ was calculated by numerical integration of
$f^\star$. (I used Simpson's rule, replacing the upper limit of the
integral with some large value.  Clearly I could have done better!)

\bigskip

$G^\star(y;\nu)$ was calculated by a nested numerical integral.  I'm
sure I could have done better!

\bigskip \bigskip

To get the MLE $\hat{\nu}$, we are in the enviable situation of
optimizing a function of one variable, $\sum_i l(\nu; y_i)$.  I used
``Brent's method'' (pulled from \emph{Numerical Recipes in C}).

\end{document}
